\begin{thebibliography}{99}
\bibitem{sixfold}
okwhynot-wtf.
\emph{Sixfold: research repository and visualisation archive}, version
\texttt{zenodo}, 18 November 2026.
\href{https://doi.org/10.5281/zenodo.24336419}{doi:10.5281/zenodo.24336419}.
Evidence commit \texttt{1c660bedf8cd3e52bad8e5d033535dcde6e06b6e}.

\bibitem{hutchinson}
J.~E. Hutchinson.
Fractals and self-similarity.
\emph{Indiana University Mathematics Journal} \textbf{30}(5) (1981), 713--747.
\href{https://doi.org/10.1512/iumj.1981.30.30055}{doi:10.1512/iumj.1981.30.30055}.

\bibitem{hata}
M. Hata.
On the structure of self-similar sets.
\emph{Japan Journal of Applied Mathematics} \textbf{2}(2) (1985), 381--414.
\href{https://doi.org/10.1007/BF03167083}{doi:10.1007/BF03167083}.

\bibitem{badgervellis}
M. Badger and V. Vellis.
H\"older parameterization of iterated function systems and a self-affine phenomenon.
\href{https://arxiv.org/abs/1910.08850v3}{arXiv:1910.08850v3}, 2020
(first posted 2019). Theorem 1.1, with a direct proof of path connectedness
without a separation assumption.

\bibitem{hochman}
M. Hochman.
On self-similar sets with overlaps and inverse theorems for entropy in $\R^d$.
\href{https://arxiv.org/abs/1503.09043v2}{arXiv:1503.09043v2}, 2017
(first posted 2015). Corollary 1.7 in the cited version.

\bibitem{garcia}
I. Garc\'ia.
Assouad dimension and local structure of self-similar sets with overlaps
in $\R^d$.
\href{https://arxiv.org/abs/1909.03148v2}{arXiv:1909.03148v2}, 2020
(first posted 2019). Corollary 1.6(1) in the cited version.

\bibitem{fhor}
J.~M. Fraser, A.~M. Henderson, E.~J. Olson and J.~C. Robinson.
On the Assouad dimension of self-similar sets with overlaps.
\emph{Advances in Mathematics} \textbf{273} (2015), 188--214.
\href{https://doi.org/10.1016/j.aim.2014.12.026}{doi:10.1016/j.aim.2014.12.026}.
Theorem 1.4 also in \href{https://arxiv.org/abs/1404.1016v2}{arXiv:1404.1016v2}.

\bibitem{olive}
M. Olive.
Effective computation of $\SO(3)$ and $\mathrm O(3)$ linear representations
symmetry classes.
\href{https://arxiv.org/abs/1709.01776v4}{arXiv:1709.01776v4}, 2019
(first posted 2017). Section 3.

\bibitem{hatcher}
A. Hatcher.
\emph{Algebraic Topology}.
Cambridge University Press, 2002. Example 1.25.
\href{https://pi.math.cornell.edu/~hatcher/AT/AT-doublepage.pdf}{Author's online edition}.

\bibitem{kunnen}
A. Kunnen and S. Schlicker.
Regular Sierpinski polyhedra.
\emph{Pi Mu Epsilon Journal} \textbf{10}(8) (1998), 607--619.
\href{https://www.pme-math.org/journal/issues/PMEJ.Vol.10.No.8.pdf}{Journal issue PDF}.

\bibitem{tsuiki}
H. Tsuiki.
Does it look square? Hexagonal bipyramids, triangular antiprismoids,
and their fractals.
\emph{Bridges 2007}, 277--286.
\href{https://archive.bridgesmathart.org/2007/bridges2007-277.html}{Proceedings record}.

\bibitem{fathauer}
R.~W. Fathauer, H. Kaczmarski and N. Duchnowski.
A fractal crystal comprised of cubes and some related fractal arrangements
of other Platonic solids.
\emph{Bridges 2008}, 289--296.
\href{https://archive.bridgesmathart.org/2008/bridges2008-289.html}{Proceedings record}.

\bibitem{syntopia}
Syntopia, with construction credited in the source to Knighty.
\emph{Octahedron.frag}, Fragmentarium, source labelled 2010;
version committed 30 July 2011.
\href{https://github.com/Syntopia/Fragmentarium/blob/14130ad5dcf8024868e7e2d34c64d911329fd83c/Fragmentarium-Source/Examples/Kaleidoscopic\%20IFS/Octahedron.frag}{Immutable source file},
commit \texttt{14130ad5dcf8024868e7e2d34c64d911329fd83c}.

\bibitem{arc}
A. Tetenov, K. Kamalutdinov and D. Vaulin.
Self-similar Jordan arcs which do not satisfy OSC.
\href{https://arxiv.org/abs/1512.00290v2}{arXiv:1512.00290v2}, 2016
(first posted 2015).

\bibitem{uniquecontact}
K. Kamalutdinov and A. Tetenov.
Even unique intersection point can break OSC: an example.
\href{https://arxiv.org/abs/1809.08595v2}{arXiv:1809.08595v2}, 2019
(first posted 2018).

\bibitem{golden}
D. Broomhead, J. Montaldi and N. Sidorov.
Golden gaskets: variations on the Sierpi\'nski sieve.
\href{https://arxiv.org/abs/math/0309304}{arXiv:math/0309304}, 2003.

\bibitem{bandt}
C. Bandt and D. Mekhontsev.
Elementary fractal geometry. 2. Carpets involving irrational rotations.
\emph{Fractal and Fractional} \textbf{6}(1) (2022), 39.
\href{https://doi.org/10.3390/fractalfract6010039}{doi:10.3390/fractalfract6010039}.

\bibitem{baranyverma}
B. B\'ar\'any and M. Verma.
Hausdorff dimension of self-similar measures and sets with common fixed-point
structure.
\emph{Ergodic Theory and Dynamical Systems}, published online 23 September 2026,
1--34.
\href{https://doi.org/10.1017/etds.2026.10342}{doi:10.1017/etds.2026.10342}.
Preprint \href{https://arxiv.org/abs/2507.05835v2}{arXiv:2507.05835v2}, 2026
(first posted 2025).
\end{thebibliography}
